\section{Original canonical crossing tails at their actual monomial leaves}

This is a source-facing consumer of the genuine interior identity empty
step.  The initial marked ideal, the actual current level and run, and the
whole selected original tails are retained.  Lower domain, testness and
cosupport are conclusions from the original canonical history.  The value
comparison is restricted explicitly to the actual order-zero branch of
that lower history; it does not assert positive-order recursive transport.
This adds the indexed-history bookkeeping needed in W\l odarczyk's Step
One and the trivial-blowup convention of Bierstone--Milman
\cite{wlodarczyk05-hironaka,bierstone-milman08-functoriality}.

\begin{theorem}[Actual canonical lower data on the original crossing split]
\label{hir:original-canonical-crossing-lower-data}
\source{wlodarczyk05-hironaka,bierstone-milman08-functoriality}
\lean{Hironaka.HironakaOriginalCrossingLowerLeaf.original_canonical_crossing_lower_data}
\leanok
\uses{hir:step-one-canonical-tail,hir:original-interior-crossing-selected-splits}
Let $n>0$, let $M$ belong to the original invariant domain in dimension
$k+1$, and assume that $M$ has maximal order.  Let $HS$ be a canonical
history for its actual Step One value built from the specific invariant
$R_k$, and let $y$ be a nonspecial point in the actual transformed
cosupport.  Put $\ell=\operatorname{level}(M,HS,y)$,
$r=\operatorname{run}_{M,\ell}(HS,y)$ and assume $|S|\leq r$.
Put $a=r-|S|$, $K=\operatorname{head}(H,a)$,
$B=\operatorname{tail}(H,a)$, $p=(BS).\operatorname{proj}(y)$,
$W=\operatorname{lowerMask}(M,\ell,K,p)$ and
$J=\operatorname{lowerMarkedM}(M,\ell,K,p)$.
Then $p\in W$, $J$ belongs to the genuine invariant domain in dimension
$k$, the entire restricted tail $(BS)|_W$ is canonical for $R_k(J)$,
and the actual point $y$ over $W$ belongs to its transformed cosupport.
\end{theorem}
\begin{proof}
The original canonical Step One history is a genuine multiple test
history both of $M$ and of $\overline M$, with identical transformed
cosupports.  Its canonicity is also the canonicity used for
$\overline M$.  The canonical-tail theorem therefore gives actual lower
data at the run split of this nonspecial cosupport point.  The true
crossing split calculation identifies its entire split record with
$(\operatorname{stage}(H,a),K,BS)$.  Substitute this whole record in
the canonical-tail conclusion.  The lower mask, marked ideal, restricted
tail and cosupport point are exactly the ones stated, rather than
separately selected representatives.  The dimension bound is
$(k+1)-1=k$.
\end{proof}

\begin{lemma}[Actual restricted interior tails at a monomial leaf]
\label{hir:original-restricted-interior-monomial-leaf}
\source{wlodarczyk05-hironaka,bierstone-milman08-functoriality}
\lean{Hironaka.HironakaOriginalCrossingLowerLeaf.restricted_interior_monomial_leaf}
\leanok
\uses{hir:original-interior-empty-test-history,hir:original-interior-empty-full-family-monomial-leaf,hir:original-crossing-restricted-empty-history,bcalc:general-value-order-zero}
Let $W$ be an actual open subset of a stage $D$, $J$ an actual marked
ideal on $W$ in the dimension-$k$ domain, and $B,S$ actual consecutive
histories from $D$.  Suppose $|S|>0$ and $(BS)|_W$ is a test history of
$J$.  At any actual point $y$ over $W$ whose nonmonomial order on this
restricted old history is zero, the specific full invariant $R_k$ on
$(B\epsilon S)|_W$ equals its value on $(BS)|_W$ under the actual
index map $\operatorname{shift}_{k,j}$, where
$j=|J.E|+|B|$.  Both points are the same underlying point $y$; membership
in the new open preimage is supplied by the identity-step projection.
\end{lemma}
\begin{proof}
The true restriction-of-concatenation isomorphism identifies the old
restricted history with $B|_W$ followed by
$S|_{B^{-1}W}$.  Its base and final identifications are the identity on
points.  Testness transfers across this isomorphism, and the genuine
order-isomorphism dictionary transfers the assumed actual order zero.
The restriction of $S$ has length $|S|>0$.  The proved same-base interior
test equivalence supplies testness after inserting the identity empty
step into this normalized history.  The actual full-family monomial
leaf theorem on this history gives the index
$|J.E|+|B|_W=|J.E|+|B|$ and the stated shift, including $k=0$.
The genuine nonmonomial-order equality supplies order zero after insertion.
Finally the whole restricted empty-history isomorphism transfers the
new value back to $(B\epsilon S)|_W$.  At order zero this uses only the
actual testness and domain: the general value is the true monomial
invariant, so no desired new canonical segment or lower-value equality
is required.  For $k>0$ this is exactly the order-zero branch of $R_k$;
for $k=0$ its value is constant.  All final identifications fix $y$.
\end{proof}

\begin{theorem}[The actual original crossing Step One value at its lower leaf]
\label{hir:original-canonical-crossing-step-one-monomial-leaf}
\source{wlodarczyk05-hironaka,bierstone-milman08-functoriality}
\lean{Hironaka.HironakaOriginalCrossingLowerLeaf.original_canonical_crossing_stepOne_leaf}
\leanok
\uses{hir:original-canonical-crossing-lower-data,hir:original-restricted-interior-monomial-leaf,hir:original-interior-crossing-lower-history,hir:original-interior-empty-step-one-dictionary}
Retain the original canonical hypotheses and literal $\ell,r,K,B,W,J,p$
above and assume $|S|>0$.  The actual lower data at $p$ exist internally.
If the nonmonomial order at their actual cosupport point on $(BS)|_W$
is zero, then the whole actual Step One value at $y$ on $H\epsilon S$
is
\[
 \operatorname{pairShift}(\operatorname{shift}_{k,|J.E|+|B|})
 \bigl(\operatorname{StepOne}(R_k,M,HS,y)\bigr).
\]
The original initial $M$ is unchanged and both Step One values use their
own actual current-level runs.  This holds at all actual levels,
including the original chosen-contact level zero.
\end{theorem}
\begin{proof}
The first theorem supplies genuine mask membership, the actual lower
domain and canonical old tail.  The resolution invariant makes that
canonical tail a test history.  The original whole history also supplies
testness of its original head $H$.  The actual crossing-run theorem
therefore gives run $r+1$ and the two entire selected split records
$(D,K,BS)$ and $(D,K,B\epsilon S)$.  Their projected point is literally
the same $p$, hence their lower mask and whole lower marked ideal are
the same $W,J$, also in the zero-contact branch.  The level dictionary
identifies both current levels with $\ell$; the special-point dictionary
makes the new point nonspecial.  Unfolding the actual Step One values
now selects these two genuine lower tails.  The mask membership supplied
by the first theorem selects the finite lower-value branch on both
histories.  The preceding restricted-leaf theorem gives the full $R_k$
shift of this true lower value.  Keeping the common first coordinate
$\ell$ gives exactly the displayed pair map.  No independently chosen
lower ideal, contact direction, tail, new canonicity or value comparison
is an input.
\end{proof}

Positive-order recursion inside this actual crossing lower tail, the full
outer general-value transport, mixed spectra and noncompact scheduling
remain further obligations.
